% Source: https://artofproblemsolving.com/community/c3h1914079_100th_post
% Stable identity: aops-topic-1914079-q02
The exponent of $3$ in $100!$ is $33+11+3+1=48$. Every divisor is uniquely $d=3^e m$ with $0\le e\le48$ and $m$ a divisor of $M=\dfrac{100!}{3^{48}}$, and choosing $d$ uniformly means choosing $e$ and $m$ independently and uniformly. If $d\equiv2\pmod9$ then $3\nmid d$, so we need $e=0$ (probability $\dfrac{1}{49}$) and $m\equiv2\pmod9$.

\textbf{$m$ is uniform on the units modulo $9$.} The exponents of $17$ and $43$ in $100!$ are $5$ and $2$. Write $m=17^j\,43^k\,r$ where $j\in\{0,\dots,5\}$, $k\in\{0,1,2\}$ and $r$ (a divisor of $\dfrac{M}{17^5\cdot43^2}$) are independent and uniform. Modulo $9$, $17\equiv-1$ and $43\equiv7$, with $7^2\equiv4$ and $7^3\equiv1$. The six units modulo $9$ are $\{1,2,4,5,7,8\}=\{\pm1,\pm4,\pm7\}$, and each is uniquely $s\cdot t$ with $s\in\{1,-1\}$ and $t\in\{1,4,7\}$. Now $17^j\equiv(-1)^j$ takes the values $1$ and $-1$ three times each, and $43^k$ takes the values $1,7,4$ once each, independently. So $17^j43^k$ is uniformly distributed over the six units. Multiplying by $r$ (a unit) permutes the units, so for every fixed $r$, and hence overall, $m$ is uniform on the six units.

Therefore the probability is
\[
\dfrac{1}{49}\cdot\dfrac{1}{6}=\boxed{\dfrac{1}{294}}.
\]
